% GENERATED FILE. Edit canonical lessons, then run npm run generate:books.
% canonical-source-hash: fnv1a64:e8cf93be406a5051
% canonical-lessons: JA-C461-takkyuu, JA-C461-kendou, JA-C461-karate, JA-C461-taisou, JA-C461-kantoku, JA-R461-review-456-461

\chapter{Japanese Sports}
\label{ch:ja-japanese-sports}
\hlchaptermodality{461}

\begin{chapteropening}
\textbf{By the end of this chapter:} \emph{I can name some Japanese sports.}

Complete the last lesson of chapter 461: I can name some Japanese sports.
\end{chapteropening}

\section[\texorpdfstring{\ja{たっきゅう} (takkyū)}{たっきゅう (takkyū)}]{\ja{たっきゅう} (takkyū) — table tennis}

\label{lesson:JA-C461-takkyuu}



\begin{quote}
Before the new one: say the Japanese for to take on a challenge, then the Japanese for to change jobs.
\end{quote}

\subsection*{You'll want to know: \ja{たっきゅう}}
\textbf{\ja{たっきゅう}} — \emph{takkyū} — \textquotedblleft{}table tennis\textquotedblright{}.

\begin{quote}
\textbf{\ja{そふは} \ja{たっきゅうが} \ja{すきです。}} — \emph{sofu wa takkyū ga suki desu} — \textquotedblleft{}My grandfather loves table tennis.\textquotedblright{}
\end{quote}

\begin{cousinweb}
\textbf{\ja{たっきゅう}} (\emph{takkyū}) is Sino-Japanese (kango): two Chinese characters read with their on'yomi, readings borrowed from Chinese. They are \emph{ta}, \textquotedblleft{}table\textquotedblright{} (Middle Chinese \emph{traewk}) and \emph{kyū}, \textquotedblleft{}ball\textquotedblright{} (\emph{gjuw}) (KANJIDIC2; Middle Chinese in Baxter's transcription).
\end{cousinweb}

\begin{grammarlens}[title={What you've built}]
One more word for describing what you lived and felt.
\end{grammarlens}

\subsection*{Your turn}
\begin{itemize}
\raggedright
  \item \emph{Say it:} \emph{takkyū}
  \item \emph{Say it:} \emph{takkyū}, once more
  \item \emph{Say it:} \emph{sofu wa takkyū ga suki desu} — My grandfather loves table tennis.
  \item \emph{Recall:} say the Japanese for to take on a challenge, then the Japanese for to change jobs, then say \emph{takkyū} again
\end{itemize}

\begin{tcolorbox}[breakable,colback=teal!4,colframe=teal!35!black,title={Before you move on}]
What does \ja{たっきゅう} mean? (\textquotedblleft{}Table tennis\textquotedblright{}.) Say it once more.
\end{tcolorbox}
\section[\texorpdfstring{\ja{けんどう} (kendō)}{けんどう (kendō)}]{\ja{けんどう} (kendō) — kendo}

\label{lesson:JA-C461-kendou}



\begin{quote}
Before the new one: say the Japanese for to change jobs, then the Japanese for table tennis.
\end{quote}

\subsection*{You'll want to know: \ja{けんどう}}
\textbf{\ja{けんどう}} — \emph{kendō} — \textquotedblleft{}kendo\textquotedblright{}.

\begin{quote}
\textbf{\ja{おとうとは} \ja{けんどうを} \ja{ならっています。}} — \emph{otōto wa kendō o naratte imasu} — \textquotedblleft{}My younger brother is learning kendo.\textquotedblright{}
\end{quote}

\begin{cousinweb}
\textbf{\ja{けんどう}} (\emph{kendō}) is Sino-Japanese (kango): two Chinese characters read with their on'yomi, readings borrowed from Chinese. They are \emph{ken}, \textquotedblleft{}sword\textquotedblright{} (Middle Chinese \emph{kjaemH}) and \emph{dō}, \textquotedblleft{}way\textquotedblright{} (\emph{dawX}) (KANJIDIC2; Middle Chinese in Baxter's transcription).
\end{cousinweb}

\begin{grammarlens}[title={What you've built}]
One more word for describing what you lived and felt.
\end{grammarlens}

\subsection*{Your turn}
\begin{itemize}
\raggedright
  \item \emph{Say it:} \emph{kendō}
  \item \emph{Say it:} \emph{kendō}, once more
  \item \emph{Say it:} \emph{otōto wa kendō o naratte imasu} — My younger brother is learning kendo.
  \item \emph{Recall:} say the Japanese for to change jobs, then the Japanese for table tennis, then say \emph{kendō} again
\end{itemize}

\begin{tcolorbox}[breakable,colback=teal!4,colframe=teal!35!black,title={Before you move on}]
What does \ja{けんどう} mean? (\textquotedblleft{}Kendo\textquotedblright{}.) Say it once more.
\end{tcolorbox}
\section[\texorpdfstring{\ja{からて} (karate)}{からて (karate)}]{\ja{からて} (karate) — karate}

\label{lesson:JA-C461-karate}



\begin{quote}
Before the new one: say the Japanese for table tennis, then the Japanese for kendo.
\end{quote}

\subsection*{You'll want to know: \ja{からて}}
\textbf{\ja{からて}} — \emph{karate} — \textquotedblleft{}karate\textquotedblright{}.

\begin{quote}
\textbf{\ja{からての} \ja{しあいで} \ja{かちました。}} — \emph{karate no shiai de kachimashita} — \textquotedblleft{}I won a karate match.\textquotedblright{}
\end{quote}

\begin{cousinweb}
\textbf{\ja{からて}} (\emph{karate}) is a native Japanese compound (wago), inherited rather than borrowed from Chinese: \emph{kara} and \emph{te}, \textquotedblleft{}hand\textquotedblright{}. JMdict gives two spellings, one writing \emph{kara} with the kanji for \textquotedblleft{}empty\textquotedblright{} and one with the kanji for \textquotedblleft{}China\textquotedblright{} (KANJIDIC2), so which \emph{kara} it holds is not settled (JMdict, KANJIDIC2).
\end{cousinweb}

\begin{grammarlens}[title={What you've built}]
One more word for describing what you lived and felt.
\end{grammarlens}

\subsection*{Your turn}
\begin{itemize}
\raggedright
  \item \emph{Say it:} \emph{karate}
  \item \emph{Say it:} \emph{karate}, once more
  \item \emph{Say it:} \emph{karate no shiai de kachimashita} — I won a karate match.
  \item \emph{Recall:} say the Japanese for table tennis, then the Japanese for kendo, then say \emph{karate} again
\end{itemize}

\begin{tcolorbox}[breakable,colback=teal!4,colframe=teal!35!black,title={Before you move on}]
What does \ja{からて} mean? (\textquotedblleft{}Karate\textquotedblright{}.) Say it once more.
\end{tcolorbox}
\section[\texorpdfstring{\ja{たいそう} (taisō)}{たいそう (taisō)}]{\ja{たいそう} (taisō) — gymnastics; physical exercises}

\label{lesson:JA-C461-taisou}



\begin{quote}
Before the new one: say the Japanese for kendo, then the Japanese for karate.
\end{quote}

\subsection*{You'll want to know: \ja{たいそう}}
\textbf{\ja{たいそう}} — \emph{taisō} — \textquotedblleft{}gymnastics; physical exercises\textquotedblright{}.

\begin{quote}
\textbf{\ja{まいあさ} \ja{たいそうを} \ja{します。}} — \emph{maiasa taisō o shimasu} — \textquotedblleft{}I do exercises every morning.\textquotedblright{}
\end{quote}

\begin{cousinweb}
\textbf{\ja{たいそう}} (\emph{taisō}) is Sino-Japanese (kango): two Chinese characters read with their on'yomi, readings borrowed from Chinese. They are \emph{tai}, \textquotedblleft{}body\textquotedblright{} (Middle Chinese \emph{thejX}) and \emph{sō}, \textquotedblleft{}maneuver\textquotedblright{} (\emph{tshaw}) (KANJIDIC2; Middle Chinese in Baxter's transcription).
\end{cousinweb}

\begin{grammarlens}[title={What you've built}]
One more word for describing what you lived and felt.
\end{grammarlens}

\subsection*{Your turn}
\begin{itemize}
\raggedright
  \item \emph{Say it:} \emph{taisō}
  \item \emph{Say it:} \emph{taisō}, once more
  \item \emph{Say it:} \emph{maiasa taisō o shimasu} — I do exercises every morning.
  \item \emph{Recall:} say the Japanese for kendo, then the Japanese for karate, then say \emph{taisō} again
\end{itemize}

\begin{tcolorbox}[breakable,colback=teal!4,colframe=teal!35!black,title={Before you move on}]
What does \ja{たいそう} mean? (\textquotedblleft{}Gymnastics; physical exercises\textquotedblright{}.) Say it once more.
\end{tcolorbox}
\section[\texorpdfstring{\ja{かんとく} (kantoku)}{かんとく (kantoku)}]{\ja{かんとく} (kantoku) — a team manager; a director}

\label{lesson:JA-C461-kantoku}



\begin{quote}
Before the new one: say the Japanese for karate, then the Japanese for gymnastics.
\end{quote}

\subsection*{You'll want to know: \ja{かんとく}}
\textbf{\ja{かんとく}} — \emph{kantoku} — \textquotedblleft{}a team manager; a director\textquotedblright{}.

\begin{quote}
\textbf{\ja{かんとくは} \ja{せんしゅに} \ja{きびしいです。}} — \emph{kantoku wa senshu ni kibishii desu} — \textquotedblleft{}The team manager is strict with the players.\textquotedblright{}
\end{quote}

\begin{cousinweb}
\textbf{\ja{かんとく}} (\emph{kantoku}) is Sino-Japanese (kango): two Chinese characters read with their on'yomi, readings borrowed from Chinese. They are \emph{kan}, \textquotedblleft{}oversee\textquotedblright{} (Middle Chinese \emph{kaem}) and \emph{toku}, \textquotedblleft{}coach\textquotedblright{} (\emph{towk}) (KANJIDIC2; Middle Chinese in Baxter's transcription).
\end{cousinweb}

\begin{grammarlens}[title={What you've built}]
One more word for describing what you lived and felt.
\end{grammarlens}

\subsection*{Your turn}
\begin{itemize}
\raggedright
  \item \emph{Say it:} \emph{kantoku}
  \item \emph{Say it:} \emph{kantoku}, once more
  \item \emph{Say it:} \emph{kantoku wa senshu ni kibishii desu} — The team manager is strict with the players.
  \item \emph{Recall:} say the Japanese for karate, then the Japanese for gymnastics, then say \emph{kantoku} again
\end{itemize}

\begin{tcolorbox}[breakable,colback=teal!4,colframe=teal!35!black,title={Before you move on}]
What does \ja{かんとく} mean? (\textquotedblleft{}A team manager; a director\textquotedblright{}.) Say it once more.
\end{tcolorbox}
\section[(dialogue)]{Looking back: from Pay and Time Off to Japanese Sports}

\label{lesson:JA-R461-review-456-461}



\begin{quote}
Say the Japanese for a salary and to earn (money).
\end{quote}

\subsection*{Your turn}
Say these aloud:

\begin{itemize}
\raggedright
  \item a salary, to earn (money), time off work, to hire, a workplace
  \item to find a job, a job interview, a curriculum vitae, to apply, recruitment
  \item a lawyer, a civil servant, an actor, a novelist, a craftsman
  \item a university, to study abroad, to do research, one's field of study, a career path
  \item a goal, to succeed, real ability, to take on a challenge, to change jobs
  \item table tennis, kendo, karate, gymnastics, a team manager
\end{itemize}

\begin{tcolorbox}[breakable,colback=teal!4,colframe=teal!35!black,title={Before you move on}]
A salary? (\textbf{\ja{きゅうりょう}}, \emph{kyūryō}.) To earn (money)? (\textbf{\ja{かせぐ}}, \emph{kasegu}.) Time off work? (\textbf{\ja{きゅうか}}, \emph{kyūka}.) To hire? (\textbf{\ja{やとう}}, \emph{yatou}.) A workplace? (\textbf{\ja{しょくば}}, \emph{shokuba}.) To find a job? (\textbf{\ja{しゅうしょくする}}, \emph{shūshoku suru}.) A job interview? (\textbf{\ja{めんせつ}}, \emph{mensetsu}.) A curriculum vitae? (\textbf{\ja{りれきしょ}}, \emph{rirekisho}.) To apply? (\textbf{\ja{おうぼする}}, \emph{ōbo suru}.) Recruitment? (\textbf{\ja{ぼしゅう}}, \emph{boshū}.) A lawyer? (\textbf{\ja{べんごし}}, \emph{bengoshi}.) A civil servant? (\textbf{\ja{こうむいん}}, \emph{kōmuin}.) An actor? (\textbf{\ja{はいゆう}}, \emph{haiyū}.) A novelist? (\textbf{\ja{さっか}}, \emph{sakka}.) A craftsman? (\textbf{\ja{しょくにん}}, \emph{shokunin}.) A university? (\textbf{\ja{だいがく}}, \emph{daigaku}.) To study abroad? (\textbf{\ja{りゅうがくする}}, \emph{ryūgaku suru}.) To do research? (\textbf{\ja{けんきゅうする}}, \emph{kenkyū suru}.) One's field of study? (\textbf{\ja{せんもん}}, \emph{senmon}.) A career path? (\textbf{\ja{しんろ}}, \emph{shinro}.) A goal? (\textbf{\ja{もくひょう}}, \emph{mokuhyō}.) To succeed? (\textbf{\ja{せいこうする}}, \emph{seikō suru}.) Real ability? (\textbf{\ja{じつりょく}}, \emph{jitsuryoku}.) To take on a challenge? (\textbf{\ja{ちょうせんする}}, \emph{chōsen suru}.) To change jobs? (\textbf{\ja{てんしょくする}}, \emph{tenshoku suru}.) Table tennis? (\textbf{\ja{たっきゅう}}, \emph{takkyū}.) Kendo? (\textbf{\ja{けんどう}}, \emph{kendō}.) Karate? (\textbf{\ja{からて}}, \emph{karate}.) Gymnastics? (\textbf{\ja{たいそう}}, \emph{taisō}.) A team manager? (\textbf{\ja{かんとく}}, \emph{kantoku}.)
\end{tcolorbox}
